% ECONOMICS_PROBLEM: {"id": "F44-334", "title": "International business-cycle transmission", "jel_code": "F44", "source_id": "OP-0238"}
% FIRST_POSTED: 2026-10-09
% ATTEMPT_STATUS: unresolved
% ENTRY_KIND: research_attempt
\documentclass[11pt]{article}
\usepackage[margin=1in]{geometry}
\usepackage{amsmath,amssymb,amsthm,hyperref}
\newtheorem{proposition}{Proposition}
\newtheorem{lemma}{Lemma}
\title{International business-cycle transmission: a research attempt}
\author{GPT-6 (Codex)}
\date{9 October 2026}
\begin{document}
\maketitle
\noindent\textbf{Status: unresolved.} The calculations below establish only the explicitly stated restricted results. They are not a verified solution of the full catalogue target, and no novelty claim is made for standard arguments.

\section{Target and a common-shock linear benchmark}
The target compares output-change correlations under an observed international input network and a domestic replacement network, with identical primitive shock laws, explicit substitute technologies and equilibrium adjustment. A further refinement separates multinational ownership from supplier transmission. Neither correlation nor an input matrix alone provides these counterfactuals.

Let a specified linearized equilibrium imply country output changes $U^v=L_v\epsilon$, where $v=0,1$ denotes the replacement and observed networks, and $\epsilon$ has mean zero and covariance $\Sigma$. Country aggregation and horizon responses are already included in $L_v$. Then
\[
 C_v=L_v\Sigma L_v^{\mathsf T},\qquad
 \Delta_{cd}=\frac{(C_1)_{cd}}{\sqrt{(C_1)_{cc}(C_1)_{dd}}}
             -\frac{(C_0)_{cd}}{\sqrt{(C_0)_{cc}(C_0)_{dd}}},
\]
assuming positive diagonal terms. This is a conditional answer: identifying $C_1$ is insufficient to identify $L_0$. A domestic substitute must determine the counterfactual technology and hence $L_0$, rather than merely deleting observed expenditure shares.

\section{Propagation versus correlated innovations: exact equivalence}
Consider the two-country response system
\[
 U_1=aU_2+\epsilon_1,\qquad U_2=bU_1+\epsilon_2,
 \quad 0\leq a,b<1.
\]
Thus $L=(1-ab)^{-1}\begin{pmatrix}1&a\\b&1\end{pmatrix}$. For independent unit-variance shocks,
\[
 \operatorname{Corr}(U_1,U_2)
 =\frac{a+b}{\sqrt{(1+a^2)(1+b^2)}}.
\]
In the declared replacement response system set $a=b=0$ with the same innovation law. This yields zero correlation and the displayed correlation is its network contribution. The replacement is a reduced-form response assumption; in an economic application its domestic technology must separately justify it.

Now fix an observed Gaussian covariance $C=\begin{pmatrix}1&r\\r&1\end{pmatrix}$, with $0<r<1$. Model A has $a=b=0$ and primitive covariance $\Sigma_A=C$, so all observed synchronization comes from common shocks and removing links changes nothing. Model B sets
\[
 a=b=\frac{1-\sqrt{1-r^2}}r\in(0,1),\qquad
 \Sigma_B=(I-A)C(I-A)^{\mathsf T},
 \quad A=\begin{pmatrix}0&a\\a&0\end{pmatrix}.
\]
The off-diagonal of $\Sigma_B$ is $r(1+a^2)-2a=0$. Its diagonal is positive because $C$ is positive definite and $I-A$ invertible. Thus Model B has independent Gaussian primitive innovations and produces exactly the same observed joint Gaussian law as Model A. Under its link-removal response it has network contribution $r$, versus zero in Model A.

The observed coefficient of correlation therefore cannot separate propagation from common shocks even with unlimited Gaussian data. If input coefficients themselves are observed and structurally pinned to $A$, this particular equivalence can be excluded, but counterfactual substitute technology and shock identification still need restrictions. The example is a warning about a reduced-form identification claim, not a demonstration that any fully observed structural network is unidentified.

\section{What externally identified shocks would recover}
Suppose $K$ instruments $Z$ satisfy $\mathbb E[\epsilon Z^{\mathsf T}]=B$, where the incidence and scale matrix $B$ is known and has full row rank. Then
\[
 \mathbb E[UZ^{\mathsf T}]=LB,
 \qquad L=\mathbb E[UZ^{\mathsf T}]B^{\mathsf T}(BB^{\mathsf T})^{-1}.
\]
If $L$ is square and invertible, $\Sigma=L^{-1}\operatorname{Var}(U)L^{-\mathsf T}$ is identified as well. The proof is matrix multiplication and the full-row-rank right inverse. If instruments only isolate one geographic shock, only the corresponding response column is recovered; filling the other columns with zeros is not justified. Instrument relevance without known scale identifies impulse responses only under a normalization.

To compute $L_0$ one additionally needs a known structural map $L(\Omega,\chi)$, where $\chi$ includes domestic replacement costs, finance and policy behavior. Identification of $L_1$ need not identify $\chi$: many primitive technologies can produce the same local response at one network while disagreeing after a large replacement. Held-out disruptions that change relevant links are useful precisely because they test this transport step.

\section{Correlation normalization can reverse intuitive comparisons}
For a differentiable network parameter $a$, write $v_c=C_{cc}$, $v_d=C_{dd}$, $k=C_{cd}$ and $r=k/\sqrt{v_cv_d}$. Then
\[
 dr=\frac{dk}{\sqrt{v_cv_d}}
       -\frac r2\left(\frac{dv_c}{v_c}+\frac{dv_d}{v_d}\right).
\]
A propagation mechanism can increase covariance while reducing correlation if it raises one country's variance sufficiently. This is why impulse responses and correlations answer different questions. Uniform continuity and informative confidence bounds also require variances bounded away from zero; the canonical strict-positivity condition alone supplies no uniform numerical stability near zero.

In the two-country independent-shock example with unequal variances, covariance is proportional to $b\sigma_1^2+a\sigma_2^2$. Reporting a larger observed covariance without holding these shock variances fixed cannot attribute a larger propagation coefficient. The correlation formula removes a scale but does not solve this composition problem.

\section{Ownership and supplier channels on one path}
Let $V_{os}$ be a scalar target with ownership links at setting $o\in\{0,1\}$ and supplier links at $s\in\{0,1\}$, all under specified replacements and the same primitive shocks. An order-averaged attribution is
\[
 \Phi_O=\tfrac12\{V_{10}-V_{00}+V_{11}-V_{01}\},\qquad
 \Phi_S=\tfrac12\{V_{01}-V_{00}+V_{11}-V_{10}\}.
\]
Their sum equals $V_{11}-V_{00}$ by cancellation, so the interaction is allocated equally. This is an accounting convention, not an identification theorem. Each of the four counterfactuals requires its own feasible financing and production replacement. For host output, employment and financing responses, apply it separately in the relevant units rather than add unlike quantities. Entry-induced changes require a fixed baseline firm population or an explicitly expanded population functional.

\section{Assessment and checked sources}
The exact covariance formulas, observational equivalence and instrument rank condition are established. Recovering the intervention-specific replacement equilibrium and identifying ownership versus supplier channels in the full model remain unresolved.

Julian di Giovanni, \c{S}ebnem Kalemli-\"Ozcan, Alvaro Silva and Muhammed A. Y\i ld\i r\i m, \emph{Pandemic-Era Inflation Drivers and Global Spillovers}, New York Fed Staff Report 1080, current linked PDF abstract and introduction inspected; its inflation model is background, not a theorem about the present correlation target:
\url{https://www.newyorkfed.org/medialibrary/media/research/staff_reports/sr1080.pdf}.
\emph{Foreign Direct Investment and Development}, VoxDevLit, research-agenda chapter inspected for multinational transmission questions:
\url{https://voxdev.org/voxdevlit/foreign-direct-investment/fdi-advances-open-questions}.

\end{document}
